\documentclass[11pt,a4paper]{article}

% ── Encoding & fonts ─────────────────────────────────────────────────────────
\usepackage[T1]{fontenc}
\usepackage[utf8]{inputenc}
\usepackage{lmodern}
\usepackage{microtype}

% ── Page geometry ─────────────────────────────────────────────────────────────
\usepackage[top=2.5cm, bottom=2.5cm, left=2.8cm, right=2.8cm]{geometry}

% ── Typography & spacing ──────────────────────────────────────────────────────
\usepackage{parskip}
\setlength{\parskip}{4pt}
\usepackage{setspace}
\setstretch{1.05}

% ── Colours ───────────────────────────────────────────────────────────────────
\usepackage[dvipsnames,table]{xcolor}
\definecolor{accent}{HTML}{1A56DB}
\definecolor{rowalt}{HTML}{F3F6FC}
\definecolor{codebackground}{HTML}{F5F5F5}
\definecolor{rulecolor}{HTML}{CCCCCC}

% ── Section headings ──────────────────────────────────────────────────────────
\usepackage{titlesec}
\titleformat{\section}{\large\bfseries\color{accent}}{}{0em}{}[\color{rulecolor}\titlerule]
\titleformat{\subsection}{\normalsize\bfseries}{}{0em}{}
\titleformat{\subsubsection}{\normalsize\itshape}{}{0em}{}
\titlespacing{\section}{0pt}{14pt}{4pt}
\titlespacing{\subsection}{0pt}{9pt}{3pt}
\titlespacing{\subsubsection}{0pt}{6pt}{2pt}

% ── Tables ────────────────────────────────────────────────────────────────────
\usepackage{booktabs}
\usepackage{tabularx}
\usepackage{array}
\renewcommand{\arraystretch}{1.22}
\newcolumntype{L}[1]{>{\raggedright\arraybackslash}p{#1}}
\newcolumntype{C}[1]{>{\centering\arraybackslash}p{#1}}

% ── Code listings ─────────────────────────────────────────────────────────────
\usepackage{listings}
\lstset{
  backgroundcolor=\color{codebackground},
  basicstyle=\ttfamily\small,
  breaklines=true,
  breakatwhitespace=true,
  frame=single,
  framerule=0.4pt,
  rulecolor=\color{rulecolor},
  xleftmargin=8pt,
  xrightmargin=8pt,
  aboveskip=6pt,
  belowskip=6pt,
  captionpos=b,
  showstringspaces=false,
}

% ── Hyperlinks ────────────────────────────────────────────────────────────────
\usepackage[hidelinks,
  pdftitle={D2 -- Enterprise Audit \& Orchestration Report},
  pdfauthor={MyPDF Engineering Team},
  colorlinks=false
]{hyperref}

% ── Header / footer ───────────────────────────────────────────────────────────
\usepackage{fancyhdr}
\pagestyle{fancy}
\fancyhf{}
\renewcommand{\headrulewidth}{0.4pt}
\fancyhead[L]{\small\textcolor{gray}{MyPDF — D2 Enterprise Audit \& Orchestration Report}}
\fancyhead[R]{\small\textcolor{gray}{28 May 2026}}
\fancyfoot[C]{\small\thepage}

% ── Misc ──────────────────────────────────────────────────────────────────────
\usepackage{enumitem}
\setlist[itemize]{itemsep=1pt, topsep=3pt, parsep=0pt}
\setlist[enumerate]{itemsep=1pt, topsep=3pt, parsep=0pt}

% ─────────────────────────────────────────────────────────────────────────────
\begin{document}

% ── Compact title block ───────────────────────────────────────────────────────
\thispagestyle{fancy}
\begin{center}
  {\LARGE\bfseries\color{accent} D2 — Enterprise Audit \& Orchestration Report}\\[6pt]
  {\normalsize MyPDF \textbullet{} Flutter (Android + Web) \textbullet{} Firebase + Supabase \textbullet{} 28 May 2026}\\[4pt]
  {\small\itshape Prepared by the MyPDF Engineering Team \textbullet{} Private — not for redistribution}
\end{center}
\vspace{2pt}
{\small This report covers four mandated audit areas: multi-agent orchestration workflow,
domain model design and state management rationale, role-based access control matrix
and Firestore security rules, and the observability infrastructure and rollback strategy
for the live OCR feature flag.}
\vspace{8pt}
\hrule
\vspace{4pt}

\tableofcontents
\vspace{6pt}
\hrule

% ─────────────────────────────────────────────────────────────────────────────
\section{Agent Workflow \& Multi-Agent Orchestration}

\subsection{The Five-Agent Model}

Development of MyPDF is assisted by a five-agent orchestration system defined in
\texttt{.claude/agents/} and governed by \texttt{docs/projectscope.md}.
Each agent is scoped to a single concern and may not edit another agent's domain.

\vspace{4pt}
\begin{tabularx}{\textwidth}{@{}L{2.6cm} L{3.4cm} L{2.6cm} L{5.0cm}@{}}
  \toprule
  \textbf{Agent} & \textbf{Concern} & \textbf{Tools} & \textbf{Edit Scope} \\
  \midrule
  \rowcolor{rowalt}
  \texttt{architect} & Layering, provider graph, plugin role boundaries &
    Read, Glob, Grep, Bash & Diagnose only \\
  \texttt{flutter\_engineer} & UI, Riverpod controllers, GoRouter, theme, \texttt{kIsWeb} &
    +Edit, Write & \texttt{lib/features/*/presentation/}, \texttt{lib/shared/}, router, \texttt{main.dart} \\
  \rowcolor{rowalt}
  \texttt{firebase\_specialist} & Firestore, Supabase, Edge Function, datasources &
    +Edit, Write & \texttt{lib/features/*/data/}, \texttt{supabase/functions/}, \texttt{firestore.rules} \\
  \texttt{qa\_engineer} & Tests, \texttt{flutter analyze}, manual QA &
    +Edit, Write & \texttt{test/} \\
  \rowcolor{rowalt}
  \texttt{security} & Auth, rules, RLS, CORS proxy, secrets, CVEs &
    Read, Glob, Grep, Bash, WebFetch & Diagnose only \\
  \bottomrule
\end{tabularx}
\vspace{2pt}

\texttt{architect} and \texttt{security} are read-only; they diagnose and report but
cannot modify source files. Only \texttt{flutter\_engineer} and
\texttt{firebase\_specialist} hold write access, each constrained to its own path namespace.

\subsection{Dispatch Protocol and Prompts Used}

All dispatches are \textbf{parallel}: the orchestrator sends a single message containing
one Task call per required agent. \textbf{Topic-map auto-routing} matches the user's intent
against a keyword table in \texttt{docs/projectscope.md}:

\vspace{4pt}
\begin{tabularx}{\textwidth}{@{}L{6.0cm} L{7.8cm}@{}}
  \toprule
  \textbf{User intent} & \textbf{Agents dispatched} \\
  \midrule
  \rowcolor{rowalt}
  Add/fix a screen, widget, or route & \texttt{flutter\_engineer} \\
  Firestore query, Supabase upload & \texttt{firebase\_specialist} \\
  \rowcolor{rowalt}
  Auth flow or PIN change & \texttt{firebase\_specialist} + \texttt{security} \\
  Security review, rules, CVE & \texttt{security} \\
  \rowcolor{rowalt}
  Tests or \texttt{flutter analyze} red & \texttt{qa\_engineer} \\
  Layering or architecture concern & \texttt{architect} \\
  \rowcolor{rowalt}
  New feature spanning data and UI & \texttt{architect} + \texttt{flutter\_engineer} + \texttt{firebase\_specialist} \\
  Pre-merge / pre-release gate & All five agents \\
  \bottomrule
\end{tabularx}
\vspace{2pt}

Explicit \texttt{audit} prefix overrides the topic map: \texttt{audit} (all five),
\texttt{audit <names>} (named subset), \texttt{audit <feature>} (all five scoped to
\texttt{lib/features/<feature>/}). The orchestrator states the routing decision in one
line before dispatching so the user can redirect.

Every prompt is \textbf{self-contained} and includes: (1) target file paths the agent
must read first, (2) the original ask verbatim, (3) which \texttt{CLAUDE.md} section
to read, and (4) the expected report shape. For example, \texttt{audit reader} dispatches
all five agents with prompts scoped to \texttt{lib/features/reader/}; the
\texttt{architect} agent must return \texttt{\#\# Violations}, \texttt{\#\# Risks}, and
\texttt{\#\# OK}; the \texttt{flutter\_engineer} must run \texttt{flutter analyze} and
report errors before claiming completion.

\subsection{Context Drift Challenges and Mitigations}

Because agents start cold on each dispatch, they risk drifting from established patterns
or creeping outside their scope. Three structural mitigations address this:

\begin{itemize}
  \item \textbf{Hard edit-scope boundaries.} Each agent's definition declares a
    \textit{What you own} list. Violations are out-of-scope and flagged by the orchestrator.
  \item \textbf{Diagnose-only agents.} \texttt{architect}, \texttt{security}, and
    \texttt{qa\_engineer} hold no write tools for production code, so even drifted output
    cannot be committed; it can only produce a report for human review.
  \item \textbf{\texttt{CLAUDE.md} as single source of truth + \texttt{flutter analyze}
    gate.} Every agent reads \texttt{CLAUDE.md} first; its explicit \textit{DO NOT} list
    serves as a checklist. A non-zero analyzer exit is an objective signal of drift.
\end{itemize}

\subsection{Handoff Management}

After all agents return, the orchestrator produces a single consolidated report grouped
by severity (\textbf{Critical / High / Medium / Info}), with each finding tagged by the
originating agent. Conflict resolution follows two tiebreakers: the \texttt{architect}
verdict prevails on structural disagreements; the \texttt{security} verdict prevails on
data-exposure disagreements.

% ─────────────────────────────────────────────────────────────────────────────
\section{Architecture \& Data}

\subsection{Clean Architecture Pattern}

Every feature in \texttt{lib/features/} follows a three-layer structure:
\texttt{data/} (SDK adapters: Firebase, Supabase, Hive, HTTP) $\rightarrow$
\texttt{domain/} (plain Dart models + repository interfaces) $\rightarrow$
\texttt{presentation/} (Riverpod controllers, screens, widgets).

The \textbf{domain layer is framework-free by rule}: zero imports of
\texttt{flutter/}, \texttt{firebase\_*}, \texttt{supabase\_*}, \texttt{hive*},
\texttt{path\_provider}, or \texttt{dart:io} are permitted inside any \texttt{domain/}
directory — enforced as a hard violation by the \texttt{architect} agent. Repository
interfaces declare return types as \texttt{Either<Failure,\,T>} from \texttt{dartz},
using \texttt{lib/core/errors/failures.dart} for the failure hierarchy
(\texttt{AuthFailure}, \texttt{ServerFailure}).

\subsection{Domain Models}

\vspace{4pt}
\begin{tabularx}{\textwidth}{@{}L{3.2cm} L{6.4cm} L{4.0cm}@{}}
  \toprule
  \textbf{Model} & \textbf{Key Fields} & \textbf{Firestore Collection} \\
  \midrule
  \rowcolor{rowalt}
  \texttt{UserModel} &
    \texttt{uid}, \texttt{name}, \texttt{email}, \texttt{role} (default \texttt{'user'}) &
    \texttt{users/\{uid\}} \\
  \texttt{BookshelfModel} &
    \texttt{id}, \texttt{name}, \texttt{ownerId}, \texttt{createdAt} &
    \texttt{bookshelves/\{shelfId\}} \\
  \rowcolor{rowalt}
  \texttt{BookModel} &
    \texttt{id}, \texttt{title}, \texttt{link}, \texttt{totalPages}, \texttt{currentPage},
    \texttt{progress}, \texttt{status}, \texttt{shelfId}, \texttt{ownerId},
    \texttt{lastReadAt?}, \texttt{author?}, \texttt{year?}, \texttt{needsOcr: bool} &
    \texttt{books/\{bookId\}} \\
  \texttt{NoteModel} &
    \texttt{id}, \texttt{bookId}, \texttt{title}, \texttt{content}, \texttt{updatedAt} &
    \texttt{notes/\{noteId\}} \\
  \bottomrule
\end{tabularx}
\vspace{2pt}

All \texttt{DateTime} fields are stored as \textbf{ISO 8601 strings} in Firestore
(not \texttt{Timestamp}), keeping the domain layer free of Firestore SDK types.
\texttt{BookModel.needsOcr} is set once at upload time by the
\texttt{\_isBitmapOnlyPdf} heuristic and is thereafter immutable (enforced by
Firestore rules): it lets the reader skip the text-extraction probe and route
directly to the Tesseract OCR pipeline for bitmap-only PDFs.

\subsection{Firestore Collection Hierarchy}

The schema uses \textbf{flat top-level collections} — \texttt{users}, \texttt{bookshelves},
\texttt{books}, \texttt{notes} — rather than sub-collections. This avoids the recursive
rule complexity that sub-collections introduce in Firestore Security Rules and enables
simple cross-shelf queries (\texttt{where('ownerId', isEqualTo: uid)}) without collection
groups. Ownership is encoded as \texttt{ownerId} fields; note ownership resolves
transitively through the parent book's \texttt{ownerId} via the \texttt{bookOwner()}
helper in Firestore rules.

Cascading deletes are handled in \texttt{FirestoreDataSource}: deleting a shelf
batch-unshelves all its books before removing the shelf doc; deleting a book
batch-deletes child notes, then returns \texttt{book.link} so the controller can
purge the Supabase object, local file cache, thumbnail cache, Hive recents entry, and
Hive OCR cache. All batches are chunked at 500 operations.

\subsection{State Management — Riverpod Justification}

\texttt{setState} is prohibited for shared state; \texttt{flutter\_riverpod} is used
throughout. The choice was driven by four requirements: (1) \textbf{compile-time safety}
— code-generated providers are typed at compile time, eliminating a class of runtime
errors; (2) \textbf{\texttt{keepAlive}} for \texttt{authStateProvider} and
\texttt{routerProvider}, preventing GoRouter from being rebuilt on widget tree rebuilds;
(3) \textbf{\texttt{StreamProvider.family}} for per-entity reactive streams (books per
shelf, notes per book), ensuring real-time UI updates without polling; and (4)
\textbf{test-overridable injection} — \texttt{featureFlagsProvider} is overridden in
\texttt{ProviderScope} so tests supply a deterministic \texttt{FeatureFlags} stub
without network calls.

\vspace{4pt}
\begin{tabularx}{\textwidth}{@{}L{4.4cm} L{3.2cm} L{6.1cm}@{}}
  \toprule
  \textbf{Provider} & \textbf{Type} & \textbf{Purpose} \\
  \midrule
  \rowcolor{rowalt}
  \texttt{authStateProvider} & \texttt{StreamProvider} & Firebase auth stream \\
  \texttt{routerProvider} & \texttt{Provider} (keepAlive) & GoRouter; watches auth + PIN state \\
  \rowcolor{rowalt}
  \texttt{shelvesProvider} & \texttt{StreamProvider} & All shelves for current user \\
  \texttt{allBooksProvider} & \texttt{StreamProvider} & All books for current user \\
  \rowcolor{rowalt}
  \texttt{booksByShelfProvider} & \texttt{StreamProvider.family} & Books in one shelf, keyed by shelfId \\
  \texttt{notesByBookProvider} & \texttt{StreamProvider.family} & Notes for one book \\
  \rowcolor{rowalt}
  \texttt{ocrPageTextProvider} & \texttt{FutureProvider.family} & OCR pipeline: cache → render → recognise → clean → cache \\
  \texttt{featureFlagsProvider} & \texttt{Provider} & Remote Config singleton (overridden in \texttt{ProviderScope}) \\
  \bottomrule
\end{tabularx}

% ─────────────────────────────────────────────────────────────────────────────
\section{Security Matrix}

\subsection{RBAC Matrix}

\vspace{4pt}
\begin{tabularx}{\textwidth}{@{}L{2.5cm} C{1.65cm} C{1.65cm} C{1.65cm} C{1.65cm} C{2.0cm} C{2.05cm}@{}}
  \toprule
  \textbf{Role} & \textbf{users} & \textbf{book\-shelves} & \textbf{books} &
  \textbf{notes} & \textbf{Supabase pdfs} & \textbf{Edge Fn} \\
  \midrule
  \rowcolor{rowalt}
  Unauthenticated & Deny & Deny & Deny & Deny &
    Deny upload; public read & Public GET \\
  User (self) &
    Read own; create; update name/email &
    Full CRUD & Full CRUD & CRUD via book &
    Upload \texttt{\{uid\}/}; public read & GET any URL \\
  \rowcolor{rowalt}
  User (other) & Deny & Deny & Deny & Deny & Deny upload; public read & GET any URL \\
  Admin & Read all & Read all & Read all & Read all & No elevated access & No elevated access \\
  \bottomrule
\end{tabularx}
\vspace{2pt}

The Edge Function \texttt{pdf-proxy} is deployed \texttt{--no-verify-jwt} (public);
SSRF defences compensate (§\ref{sec:ssrf}). The \texttt{pdfs} bucket is
public-read by design (URLs are unguessable millis-timestamp paths). The
\texttt{admin} role is read-only by rule and cannot be self-assigned from the client.

\subsection{Firestore Security Rules}

\textbf{Source:} \texttt{firestore.rules}

Five helper functions encapsulate repeated logic: \texttt{isSignedIn()},
\texttt{isOwner(uid)}, \texttt{isAdmin()} (resolves \texttt{users/\{uid\}.role}),
\texttt{timestampValid(field)} (accepts Firestore \texttt{timestamp} within ±5 min
\emph{or} non-empty ISO 8601 string), and \texttt{bookOwner(bookId)} (fetches the
parent book's \texttt{ownerId} for transitive note ownership).

\textbf{users/\{uid\}:} Owner-only read/write. \texttt{create} locks \texttt{role} to
\texttt{'user'} (prevents client-side privilege escalation). \texttt{update} restricts
mutable fields to \texttt{name} and \texttt{email} via
\texttt{changedFields().hasOnly(['name','email'])}. Delete is denied.

\textbf{bookshelves/\{shelfId\}:} \texttt{create} binds \texttt{ownerId} to
\texttt{request.auth.uid}. \texttt{update} allows only the \texttt{name} field;
\texttt{ownerId} and \texttt{createdAt} are permanently locked.

\textbf{books/\{bookId\}:} \texttt{create} enforces type correctness and restricts
\texttt{status} to \texttt{['reading', 'on\_hold', 'finished']}. \texttt{update} allows
only six mutable fields: \texttt{title}, \texttt{currentPage}, \texttt{progress},
\texttt{status}, \texttt{lastReadAt}, \texttt{shelfId}. Fields \texttt{link},
\texttt{totalPages}, \texttt{needsOcr}, \texttt{author}, and \texttt{year} are immutable
after creation — a client cannot rewrite the PDF URL.

\textbf{notes/\{noteId\}:} Ownership resolves transitively via
\texttt{bookOwner(resource.data.bookId)}. Even knowledge of a note's document ID
grants no access without ownership of the parent book. Mutable fields:
\texttt{title}, \texttt{content}, \texttt{updatedAt}.

\textbf{Default deny:} A catch-all \texttt{match /\{document=**\}} denies all
reads and writes for any unmatched path.

\subsection{Supabase Storage}

The \texttt{pdfs} bucket uses path pattern \texttt{\{uid\}/\{millis\}.pdf}. RLS
policies restrict uploads to the caller's own \texttt{\{uid\}/} prefix. Deletes are
owner-enforced by the Supabase auth context. Reads are intentionally public via
\texttt{getPublicUrl()}.

\subsection{Edge Function SSRF Defences}
\label{sec:ssrf}

\textbf{Source:} \texttt{supabase/functions/pdf-proxy/index.ts}

\begin{itemize}
  \item \textbf{Private IP + hostname rejection.} DNS A/AAAA records are checked against
    all private/reserved IPv4 ranges (\texttt{10/8}, \texttt{172.16/12},
    \texttt{192.168/16}, \texttt{127/8}, \texttt{169.254/16}) and IPv6 private ranges.
    \texttt{localhost}, \texttt{*.local}, \texttt{*.internal}, and \texttt{metadata.*} are
    blocked by name before resolution.
  \item \textbf{Redirect resistance.} \texttt{fetch} uses \texttt{redirect: manual},
    preventing redirect-based SSRF through public-to-private redirects.
  \item \textbf{Size cap + timeout.} Declared \texttt{Content-Length} and streamed bytes
    are both capped at 50\,MB. A 120-second fetch timeout prevents upstream hangs.
  \item \textbf{Forced content type.} The response \texttt{Content-Type} is always
    overwritten to \texttt{application/pdf}, preventing HTML injection through the proxy.
\end{itemize}

\subsection{App-Level PIN and Biometric}

A mandatory 6-digit PIN gates every cold start
(\texttt{lib/core/local/app\_pin\_service.dart}). The PIN is hashed using
\textbf{SHA-256-crypt} (\texttt{\$5\$<salt>\$<hash>}) via \texttt{package:crypt}; the
raw PIN is never stored. Session state lives in \texttt{AppPinSession} (in-memory only,
reset on logout). Biometric (\texttt{BiometricAuthService}) is an optional acceleration
layer at the login and PIN entry screens — never a sole credential; web always returns
\texttt{false} from \texttt{isDeviceSupported()}.

\subsection{Known Security Debts}

Documented in \texttt{CLAUDE.md} as deliberately deferred for demo scope:
\texttt{pdf.js} loaded without SRI; no CSP/COEP/COOP headers (COEP absence
degrades Tesseract.js SIMD WASM silently); Hive \texttt{app\_prefs} unencrypted;
\texttt{tesseract\_ocr 0.5.0} unmaintained since 2023.

% ─────────────────────────────────────────────────────────────────────────────
\section{Observability \& Rollback}

\subsection{Structured Logging (AppLogger)}

\textbf{Source:} \texttt{lib/core/logging/app\_logger.dart}

All logging is centralised through the static \texttt{AppLogger} class (wraps
\texttt{package:logger}):

\begin{lstlisting}[language=Dart]
AppLogger.debug(tag, message, {error})   // Level.debug — suppressed in release
AppLogger.info(tag, message)
AppLogger.warn(tag, message, {error, stackTrace})
AppLogger.error(tag, message, {error, stackTrace})  // mirrors to Crashlytics
\end{lstlisting}

Tags use ALL\_CAPS\_SNAKE (\texttt{AUTH}, \texttt{OCR}, \texttt{TTS}, \texttt{LIBRARY})
for dashboard filtering. Debug builds use \texttt{PrettyPrinter} (8-frame stacks,
colored); release builds use \texttt{SimplePrinter} (timestamps, \texttt{Level.info}
threshold).

\subsection{Crashlytics Implementation}

\textbf{Sources:} \texttt{lib/main.dart}, \texttt{lib/core/logging/app\_logger.dart}

All four instrumentation sites are wrapped in \texttt{if (!kIsWeb)};
collection is further gated by
\texttt{setCrashlyticsCollectionEnabled(!kDebugMode)}.

\vspace{4pt}
\begin{tabularx}{\textwidth}{@{}C{0.4cm} L{3.2cm} L{3.2cm} L{6.4cm}@{}}
  \toprule
  \textbf{\#} & \textbf{Source} & \textbf{API call} & \textbf{What it captures} \\
  \midrule
  \rowcolor{rowalt}
  1 & \texttt{main.dart} & \texttt{recordFlutterFatalError} &
    Flutter widget build/render errors (\texttt{FlutterError.onError}) \\
  2 & \texttt{main.dart} & \texttt{recordError(fatal:true)} &
    Platform-thread and Dart isolate errors (\texttt{PlatformDispatcher.onError}) \\
  \rowcolor{rowalt}
  3 & \texttt{main.dart} & \texttt{recordError(fatal:true)} &
    Fallback — errors not caught by sites 1 or 2 (\texttt{runZonedGuarded}) \\
  4 & \texttt{app\_logger.dart} & \texttt{recordError} + \texttt{reason} &
    Non-fatal handled errors via \texttt{AppLogger.error()} \\
  \bottomrule
\end{tabularx}
\vspace{2pt}

Site 4 means any \texttt{AppLogger.error()} call with a non-null \texttt{error} object
is automatically mirrored to Crashlytics with the tag and message as the \texttt{reason}
field, surfacing non-fatal handled errors alongside fatal crashes.

\subsection{Feature Flag System}

\textbf{Source:} \texttt{lib/core/config/feature\_flags.dart}

\texttt{FeatureFlags} wraps \texttt{FirebaseRemoteConfig} with typed getters. It is
constructed once in \texttt{main()} and injected via
\texttt{featureFlagsProvider.overrideWithValue(featureFlags)}.
Initialization is \textbf{failure-safe}: a try/catch logs the error and falls through to
defaults so the app never blocks on Remote Config (fetch timeout: 10 s; cache TTL: 1 h).

\vspace{4pt}
\begin{tabularx}{\textwidth}{@{}L{4.2cm} C{1.8cm} L{7.8cm}@{}}
  \toprule
  \textbf{Remote Config key} & \textbf{Default} & \textbf{Effect} \\
  \midrule
  \rowcolor{rowalt}
  \texttt{ocr\_fallback\_enabled} & \texttt{true} &
    Master switch for the Tesseract OCR pipeline. \texttt{false} = embedded text only;
    scanned PDFs return empty TTS output rather than invoking Tesseract. \\
  \bottomrule
\end{tabularx}

\subsection{Rollback Plan — \texttt{ocr\_fallback\_enabled}}

The OCR pipeline executes WASM/FFI code on device, performs large memory allocations
for page rasterisation, and can silently fail on low-RAM devices. The flag provides a
\textbf{zero-redeploy kill switch}:

\begin{enumerate}
  \item \textbf{Trigger.} Crashlytics OCR-related crash spike, or user reports of reader
    freezing on scanned documents.
  \item \textbf{Action.} Firebase Console $\rightarrow$ Remote Config $\rightarrow$
    \texttt{ocr\_fallback\_enabled} $\rightarrow$ \texttt{false} $\rightarrow$ publish.
  \item \textbf{Propagation.} Clients pick up the change within the 1-hour TTL; no
    client redeploy is required.
  \item \textbf{Effect.} \texttt{ReadingScreen} skips the OCR branch. TTS falls back to
    the embedded text layer only — no crash, no freeze.
  \item \textbf{Recovery.} After root-cause fix is deployed, flip back to \texttt{true};
    clients update within 1 hour.
\end{enumerate}

More broadly, \texttt{CLAUDE.md} classifies all engineering actions into three
reversibility tiers: \textbf{R0} (irreversible — stop and ask first),
\textbf{R1} (costly to reverse — proceed with stated justification), and
\textbf{R2} (easily reversed — just do it). This ensures the most dangerous operations
require explicit human approval while routine development remains friction-free.

\end{document}
